% ============================================================
%  Lernzettel Template (my_dhbw Stil, analysis.tex)
%  Eine Spalte, mit TOC, accent-Theme.
%  Renderer: pdflatex (zweimal)
% ============================================================
\documentclass[11pt, a4paper]{article}

\usepackage[ngerman]{babel}
\usepackage{fontspec}
\setmainfont[
  Path=/usr/share/fonts/TTF/,
  Extension=.ttf,
  BoldFont=DejaVuSans-Bold,
  ItalicFont=DejaVuSans-Oblique,
  BoldItalicFont=DejaVuSans-BoldOblique
]{DejaVuSans}
\setsansfont[
  Path=/usr/share/fonts/TTF/,
  Extension=.ttf,
  BoldFont=DejaVuSans-Bold,
  ItalicFont=DejaVuSans-Oblique,
  BoldItalicFont=DejaVuSans-BoldOblique
]{DejaVuSans}
\setmonofont[Path=/usr/share/fonts/TTF/,Extension=.ttf]{DejaVuSansMono}
\usepackage[top=2cm, bottom=2cm, left=2cm, right=2cm, footskip=1cm]{geometry}
\usepackage{amsmath, amssymb}
\usepackage{xcolor}
\usepackage{titlesec}
\usepackage{titletoc}
\usepackage{enumitem}
\usepackage{fancyhdr}
\usepackage{lastpage}
\usepackage{microtype}
\usepackage{parskip}

\usepackage{booktabs}
\usepackage{array}
\usepackage{longtable}
\usepackage{tabularx}
\usepackage{ltablex}
\keepXColumns
\usepackage{colortbl}
\usepackage[most,breakable]{tcolorbox}
\usepackage{listings}
\usepackage[normalem]{ulem}
\usepackage{pdflscape}
\usepackage[colorlinks=true, linkcolor=accent, urlcolor=accent]{hyperref}

\renewcommand{\familydefault}{\sfdefault}

% ── Theme (my_dhbw accent) ────────────────────────────────────
\definecolor{accent}{RGB}{0, 60, 113}
\definecolor{primaryblue}{RGB}{0, 60, 113}
\definecolor{codebg}{RGB}{245, 245, 245}
\definecolor{figurebg}{RGB}{232, 241, 255}
\definecolor{tablehintbg}{RGB}{240, 255, 240}
\definecolor{solutionbg}{RGB}{242, 255, 242}
\definecolor{rulegray}{RGB}{180, 180, 180}
\definecolor{tableheadbg}{RGB}{0, 60, 113}

% ── Header/Footer ─────────────────────────────────────────────
\pagestyle{fancy}
\fancyhf{}
\renewcommand{\headrulewidth}{0pt}
\renewcommand{\footrulewidth}{0.4pt}
\renewcommand{\sectionmark}[1]{\markboth{#1}{}}
\fancyfoot[L]{\footnotesize\color{accent}\nouppercase{\leftmark}}
\fancyfoot[R]{\footnotesize\color{accent}\thepage\,/\,\pageref{LastPage}}
\fancypagestyle{plain}{%
  \fancyhf{}%
  \renewcommand{\headrulewidth}{0pt}%
  \renewcommand{\footrulewidth}{0.4pt}%
  \fancyfoot[L]{\footnotesize\color{accent}\nouppercase{\leftmark}}%
  \fancyfoot[R]{\footnotesize\color{accent}\thepage\,/\,\pageref{LastPage}}%
}

% ── Typografie ────────────────────────────────────────────────
\titleformat{\section}
    {\color{accent}\large\bfseries}{}{0pt}{}
    [\vspace{-4pt}\color{accent}\rule{\linewidth}{2pt}]
\titleformat{\subsection}
    {\normalsize\bfseries\color{accent!85!black}}{}{0pt}{}
    [\vspace{-3pt}\color{accent!40}\rule{\linewidth}{0.4pt}]
\titleformat{\subsubsection}
    {\small\bfseries\color{accent!70!black}}{}{0pt}{}{}
\titlespacing*{\section}{0pt}{12pt}{4pt}
\titlespacing*{\subsection}{0pt}{8pt}{2pt}
\titlespacing*{\subsubsection}{0pt}{6pt}{1pt}

\setlength{\parindent}{0pt}
\setlength{\parskip}{4pt}
\renewcommand{\arraystretch}{1.25}
\setlist[itemize]{noitemsep, topsep=3pt, parsep=0pt, leftmargin=1.4em}
\setlist[enumerate]{noitemsep, topsep=3pt, parsep=0pt, leftmargin=1.6em}
\setlist[itemize,2]{label=\textendash, leftmargin=1.4em}

% ── Boxen ──────────────────────────────────────────────────────
\newtcolorbox{figurebox}[1]{
  colback=figurebg, colframe=accent!50,
  title={Abbildung: #1}, fonttitle=\small\bfseries,
  breakable, before skip=6pt, after skip=6pt
}
\newtcolorbox{tablehintbox}[1]{
  colback=tablehintbg, colframe=green!50!black!50,
  title={Tabelle: #1}, fonttitle=\small\bfseries,
  breakable, before skip=6pt, after skip=6pt
}
\newtcolorbox{solutionbox}[1]{
  colback=solutionbg, colframe=green!60!black!60,
  title={#1}, fonttitle=\small\bfseries,
  breakable, before skip=6pt, after skip=6pt
}

\lstset{
  basicstyle=\ttfamily\small,
  backgroundcolor=\color{codebg},
  breaklines=true,
  frame=single, framerule=0pt,
  xleftmargin=4pt, xrightmargin=4pt,
  aboveskip=6pt, belowskip=6pt,
  keepspaces=true
}

% ── TOC-Look ──────────────────────────────────────────────────
\renewcommand{\contentsname}{\color{accent}Inhalt}

\providecommand{\nichttitle}{Lernzettel}
\providecommand{\nichtsubtitle}{}

\renewcommand{\nichttitle}{Betriebssysteme}
\renewcommand{\nichtsubtitle}{Lernzettel}


\begin{document}

\begin{center}
{\bfseries\Large\color{accent}\nichttitle}\\[2pt]
\color{accent}\rule{0.35\linewidth}{1pt}\color{black}
\end{center}
\vspace{2pt}

{\small\tableofcontents}
\newpage

% ── BODY ──────────────────────────────────────────────────────

\section{Grundlagen}

\textbf{Betriebssystem (BS)}: Systemsoftware als Schnittstelle zwischen Hardware und Software; steuert/überwacht Programmabwicklung, regelt Rechnerbetrieb und stellt Benutzern/Anwendungen elementare Dienste bereit.

\begin{itemize}
  \item Kern = \textbf{Kernel}; darüber Benutzer- und Systemprozesse
\end{itemize}

\textbf{Aufgaben}: Prozess-, Speicher-, Geräte-, Datei-, Energieverwaltung.


\textbf{Programmarten}:

\begin{itemize}
  \item \textbf{Anwendungsprogramm}: löst Probleme der Benutzer
  \item \textbf{Systemprogramm}: verwaltet den Rechnerbetrieb
  \item \textbf{Betriebssystem}: stellt Benutzern/Anwendungen elementare Dienste bereit
\end{itemize}

\textbf{Benutzerschnittstellen}:

\begin{itemize}
  \item \textbf{GUI}: grafisch
  \item \textbf{CLI}: Kommandozeile
  \item \textbf{TUI}: zeichenorientiert
  \item \textbf{VUI}: Sprache
  \item \textbf{NUI}: natürlich
\end{itemize}


\section{Geschichte (Meilensteine)}


{\small
\begin{tabularx}{\linewidth}{XX}
\hline
\rowcolor{tableheadbg}
{\color{white}\textbf{Jahr}} & {\color{white}\textbf{Ereignis}} \\[2pt]
\hline
\endhead
\hline
\endfoot
70 v.Chr. & Mechanismus von Antikythera \\
1623 & Schickard: erste mechan. Rechenmaschine \\
1642 & Pascal: „Pascaline" (Addition 6-stellig) \\
1801 & Jacquard: lochkartengesteuerter Webstuhl \\
1832 & Babbage: Difference/Analytical Engine; Lovelace: erstes Programm (Bernoulli) \\
1887 & Hollerith: Lochkarten-Tabulator → IBM \\
1905 & Fleming: Vakuumdiode \\
1938/41 & Zuse Z1/Z3: erster binär programmgesteuerter Rechner \\
1944 & Harvard Mark I (Aiken) \\
1946 & ENIAC: erster Röhrenrechner \\
1947 & Williams-Kilburn-Röhre; erster Transistor \\
1958/59 & Kilby/Noyce: integrierter Schaltkreis \\
1964 & IBM OS/360 \\
1968 & Unix (Ritchie/Thompson); Engelbart: Maus \\
1969 & ARPANET; Apollo 11 (Hamilton) \\
1971 & Intel 4004 (erster Mikroprozessor) \\
1976 & Apple I; IP \\
1981 & IBM 5150 PC \\
1985 & Windows 1.0 \\
1987/90 & Minix (Tanenbaum) → Linux (Torvalds) \\
1995 & Windows 95 (eigenständig, 32-Bit) \\
2007 & iPhone \\
2023 & Frontier-Supercomputer: 1,194 EFLOPS \\
\end{tabularx}
}


\textbf{BS-Fehler-Fallbeispiele}:

\begin{itemize}
  \item \textbf{Ariane 5}: 64-Bit Float → 16-Bit Int Overflow
  \item \textbf{Mars Climate Orbiter}: SI vs. imperial (Faktor 4,45)
  \item \textbf{F-16}: negative Breitengrade nicht behandelt
  \item \textbf{Spectre/Meltdown}: CPU-Designfehler
\end{itemize}


\section{Rechnerarchitektur}

\textbf{Hauptkomponenten}: CPU, Speicher (Haupt + Sekundär), E/A-Geräte, Busse.


\textbf{Speicherhierarchie} (Geschwindigkeit ↓, Größe ↑): Register → Cache → Hauptspeicher (RAM/ROM) → Sekundärspeicher.


\textbf{CPU}: Central Processing Unit — führt Befehle aus, steuert Abläufe.

\begin{itemize}
  \item \textbf{Steuerwerk} = Befehlsprozessor
  \item \textbf{ALU/Rechenwerk} = Datenprozessor (mit Status/Flags, Akkumulator)
  \item Register: IR, MBR, PC, MAR
  \item \textbf{Systembus}: Adress-, Daten-, Steuerbus
\end{itemize}

\textbf{Registerbreite}: Anzahl gleichzeitig speicherbarer Bits (4/8/16/32/64). 64-Bit-OS adressiert 2⁶⁴ Speichereinheiten.



\section{Modi und Systemaufruf}


{\small
\begin{tabularx}{\linewidth}{XXX}
\hline
\rowcolor{tableheadbg}
{\color{white}\textbf{Eigenschaft}} & {\color{white}\textbf{User-Mode}} & {\color{white}\textbf{Kernel-Mode}} \\[2pt]
\hline
\endhead
\hline
\endfoot
Privileg & nicht-privilegiert & privilegiert \\
Nutzer & Anwendungen & BS-Code \\
Befehlssatz & eingeschränkt & beliebig \\
Speicher & kein Kernel-Zugriff & gesamter Speicher + alle Betriebsmittel \\
\end{tabularx}
}


\textbf{Modusübergänge}:

\begin{itemize}
  \item Kernel→User: unproblematisch, durch BS bei CPU-Zuteilung
  \item User→Kernel: nur indirekt über Systemaufruf
\end{itemize}

\textbf{Systemaufruf (Syscall)}: Aufruf einer BS-Funktion durch User-Mode-Prozess; nur im Kernel-Mode ausführbar.


Ablauf:

\begin{enumerate}
  \item User-Prozess tätigt Syscall
  \item Umschalten in Kernel-Mode
  \item Sicherheitsprüfung (Berechtigung)
  \item Funktion ausführen oder ablehnen
  \item Rückkehr in User-Mode mit Rückgabewert
\end{enumerate}

Bsp. Dateizugriff: \texttt{open}, \texttt{read}, \texttt{write}, \texttt{close}.



\section{Betriebsmittel}

\textbf{Betriebsmittel (Ressource)}: Hard- oder Software-Ressource des Rechners.



{\small
\begin{tabularx}{\linewidth}{XXX}
\hline
\rowcolor{tableheadbg}
{\color{white}\textbf{Art}} & {\color{white}\textbf{Definition}} & {\color{white}\textbf{Beispiel}} \\[2pt]
\hline
\endhead
\hline
\endfoot
Entziehbar & jederzeit ohne Schaden entziehbar & CPU \\
Nicht entziehbar & muss bis Beendigung bleiben & Drucker \\
Exklusiv & max. 1 Prozess gleichzeitig & — \\
Gemeinsam & quasi-gleichzeitig nutzbar & Festplatte \\
\end{tabularx}
}


\textbf{BS-Kategorien}: Mainframes (z/OS), Server, PC/Laptop, Echtzeit, Embedded, Tablets, Handhelds, Smartcards.



\section{Prozessverwaltung}

\textbf{Prozess}: Programm zur Laufzeit / Ablaufumgebung; dynamische Folge von Aktionen mit Zustandsänderungen. Adressraum: Userspace + Kernelspace.


\textbf{Programm}: statische Verfahrensvorschrift.


\textbf{Programm → Prozess}: Compiler übersetzt → BS lädt in Adressraum.



\subsection{Prozesserzeugung}
\begin{enumerate}
  \item Code/Daten in Speicher laden
  \item Eindeutige PID zuordnen
  \item PCB in Prozesstabelle anlegen
  \item Syscalls: Unix \texttt{fork}, Windows \texttt{CreateProcess}
\end{enumerate}


\subsection{Prozesskenndaten}
\begin{itemize}
  \item \textbf{PID}: ganze Zahl ≥ 0, eindeutig
  \item \textbf{PPID}: Elternprozess
  \item UID/GID, EUID/EGID; PRI; Zustand; verbrauchte Rechenzeit
\end{itemize}


{\small
\begin{tabularx}{\linewidth}{XXX}
\hline
\rowcolor{tableheadbg}
{\color{white}\textbf{BS}} & {\color{white}\textbf{PID 0}} & {\color{white}\textbf{PID 1}} \\[2pt]
\hline
\endhead
\hline
\endfoot
Unix & Idle & Init \\
Windows & nicht vergeben & nicht vergeben \\
\end{tabularx}
}



\subsection{Prozesskontext \& PCB}

\textbf{Prozesskontext}: alle Infos, die ein Prozess während Ausführung braucht.


\textbf{Prozesstabelle}: enthält PCBs aller Prozesse.


\textbf{PCB-Inhalt}:

\begin{itemize}
  \item Prozesszustand, Programmzähler, PID
  \item Initiale + dynamische Priorität
  \item Verbrauchte CPU-Zeit
  \item Zugeordnete Betriebsmittel
  \item Aktuelle Registerinhalte
  \item Hardwarekontext: CPU-Register (PC, IR, Stack, Flags), MMU-Register, Stack/Heap/Code, Zugriffsinfos
\end{itemize}


\subsection{Kontextwechsel}
\begin{enumerate}
  \item Hardwarekontext laufender Prozess sichern (PCB\_A)
  \item Hardwarekontext neuer Prozess laden (PCB\_B)
  \item Bei Rückwechsel: umgekehrt
\end{enumerate}


\subsection{Prozesszyklus}

Zustände: \textbf{Bereit (ready)}, \textbf{Aktiv (running)}, \textbf{Blockiert (waiting)}, \textbf{Beendet}.


Übergänge:

\begin{enumerate}
  \item Aktivieren: BS wählt Prozess (Bereit→Aktiv)
  \item Deaktivieren/Vorrangunterbrechung (Aktiv→Bereit)
  \item Blockieren: wartet auf I/O/Ressource (Aktiv→Blockiert)
  \item Deblockieren: Ressource verfügbar (Blockiert→Bereit)
  \item Terminieren
\end{enumerate}

\textbf{Sonderzustände}:

\begin{itemize}
  \item \textbf{Zombie}: beendeter Prozess, noch in Tabelle; kein CPU-/Speicherverbrauch außer Eintrag; Eltern muss austragen
  \item \textbf{Dämon (Daemon)}: Hintergrundprozess unabhängig vom Terminal
\end{itemize}


\subsection{Thread}

\textbf{Thread}: kleinste Ausführungseinheit, „Prozess im Prozess", kontinuierliche Anweisungsfolge.

\begin{itemize}
  \item Prozess = mind. 1 Thread (Main Thread); endet, wenn alle Threads beendet
  \item Threads teilen Code/Daten/Heap/Ressourcen des Prozesses
\end{itemize}


{\small
\begin{tabularx}{\linewidth}{XXX}
\hline
\rowcolor{tableheadbg}
{\color{white}\textbf{Aspekt}} & {\color{white}\textbf{Prozess}} & {\color{white}\textbf{Thread}} \\[2pt]
\hline
\endhead
\hline
\endfoot
Gewicht & schwer & leicht \\
Adressraum & eigener & gemeinsam \\
Speicherschutz & ja & nein \\
Erzeugung & Namensraum dupliziert & minimaler Aufwand \\
Kontextwechsel & aufwendig & schnell \\
\end{tabularx}
}


\textbf{Vorteile Thread}: schnellerer Kontextwechsel, Zugriff auf alle Prozess-Ressourcen, Modularisierung.

\textbf{Nachteile}: Synchronisation nötig, kein Speicherschutz (Absturz reißt alle mit), spezielle Programmierkenntnisse.


\textbf{Threadzustände}: Laufend, Bereit, Wartend.


\textbf{Wechsel-Auslöser}: Zeitscheibe abgelaufen, warten auf Objekt, höher priorisierter Thread bereit.


\textbf{Threadprioritäten}:


{\small
\begin{tabularx}{\linewidth}{XX}
\hline
\rowcolor{tableheadbg}
{\color{white}\textbf{System}} & {\color{white}\textbf{Bereich}} \\[2pt]
\hline
\endhead
\hline
\endfoot
Windows & 0: System; 1–15: Timesharing; 16–31: Realtime \\
Linux & 0–99: Realtime; 100–139: Timesharing \\
\end{tabularx}
}


\textbf{Garbage Collector}: niedrige Priorität, läuft bei Leerlauf/Ressourcenknappheit.



\section{Scheduling}

\textbf{Scheduling}: Aufteilung der CPU-Zeit auf Prozesse.


\textbf{Komponenten}:

\begin{itemize}
  \item \textbf{Prozessmanager} = Scheduler + Dispatcher
  \item \textbf{Scheduler}: entscheidet, welcher Prozess die CPU bekommt
  \item \textbf{Dispatcher}: führt Kontextwechsel aus
\end{itemize}

\textbf{Job}: parallel/nebenläufig auszuführender Auftrag.



\subsection{Scheduling-Arten (Horizont)}
\begin{enumerate}
  \item Kurzfristig (CPU-Vergabe)
  \item Mittelfristig (Speicher)
  \item Langfristig (eingehende Aufgaben)
\end{enumerate}


\subsection{Verfahrensklassen}

{\small
\begin{tabularx}{\linewidth}{XX}
\hline
\rowcolor{tableheadbg}
{\color{white}\textbf{Variante}} & {\color{white}\textbf{Bedeutung}} \\[2pt]
\hline
\endhead
\hline
\endfoot
\textbf{Non-preemptive} & nicht-verdrängend; CPU nur bei Block/Freigabe/Ende \\
\textbf{Preemptive} & verdrängend; CPU jederzeit entziehbar \\
\end{tabularx}
}



\subsection{Scheduling-Ziele}
\begin{enumerate}
  \item Fairness — Mindestzuteilung
  \item Effizienz — volle CPU-Auslastung
  \item Min. Antwortzeit
  \item Min. Wartezeit
  \item Max. Durchsatz
  \item Min. Durchlaufzeit (Verweilzeit)
  \item Vorhersehbarkeit
\end{enumerate}


{\small
\begin{tabularx}{\linewidth}{XX}
\hline
\rowcolor{tableheadbg}
{\color{white}\textbf{Begriff}} & {\color{white}\textbf{Definition}} \\[2pt]
\hline
\endhead
\hline
\endfoot
Antwortzeit & Wartezeit Anwender auf Bearbeitung \\
Wartezeit & Summe aller Wartezeiträume \\
Durchsatz & Prozesse/Zeiteinheit \\
CPU-Auslastung & \% der Gesamtkapazität \\
Durchlaufzeit & Bedienzeit + Wartezeiten \\
Bedienzeit & Zeit, in der Prozess CPU hält \\
\end{tabularx}
}



\subsection{Batch-Scheduling}
\begin{itemize}
  \item \textbf{FCFS} (First Come First Serve): Reihenfolge des Eintreffens
  \item \textbf{SJF/SPF/SPN} (Shortest Job First): kürzester Job zuerst
  \item \textbf{SRTN} (Shortest Remaining Time Next): preemptive SJF; bevorzugt neue kurze Jobs
\end{itemize}

\begin{quote}
SJF/SRTN: Bedienzeit praktisch unbekannt; \textbf{Verhungern (Starvation)} möglich.
\end{quote}


\subsection{Dialog-Scheduling}
\begin{itemize}
  \item \textbf{Round Robin (RR)}: FCFS + Zeitscheibe (Quantum); alle gleich; abgelaufen → an Warteschlangen-Ende
  \item \textbf{Priority Scheduling (PS)}: höchste Priorität zuerst; gleiche Prio → eine Klasse (Multilevel)
  \item \textbf{Fair-Share-Scheduling (FSS)}: jeder von n Prozessen → 1/n CPU-Zeit; schlechtestes Verhältnis verbraucht/zustehend zuerst
  \item \textbf{Lottery Scheduling}: zufällig; n Verlosungen/s, Gewinner erhält CPU für m ms
\end{itemize}

\textbf{Quantum-Wahl}:

\begin{itemize}
  \item zu kurz → Overhead
  \item zu lang → Verzögerungen
  \item typisch \textbf{10–200 ms}; höhere Taktrate → kürzeres Quantum
  \item statisch oder dynamisch (E/A-intensiv: höher; rechenintensiv: kürzer)
\end{itemize}

\textbf{Priorität}: statisch (Voreinstellung) + dynamisch (Anpassung) — meist kombiniert.



\subsection{Realtime-Scheduling}
\begin{itemize}
  \item \textbf{EDF/MDF}: Earliest/Minimal Deadline First — kleinste nächste Deadline
  \item \textbf{Polled Loop}: zyklische Abfrage
  \item \textbf{Interrupt-gesteuert}: ISR auf Ereignis
\end{itemize}


{\small
\begin{tabularx}{\linewidth}{XX}
\hline
\rowcolor{tableheadbg}
{\color{white}\textbf{Typ}} & {\color{white}\textbf{Bedeutung}} \\[2pt]
\hline
\endhead
\hline
\endfoot
Hard Real Time & feste Zeitschranke; Verspätung = schwerer Fehler \\
Soft Real Time & weiche Schranke; Verspätung = lästig \\
\end{tabularx}
}



\subsection{Multi-Level-Scheduling}


{\small
\begin{tabularx}{\linewidth}{XX}
\hline
\rowcolor{tableheadbg}
{\color{white}\textbf{Prio}} & {\color{white}\textbf{Klasse}} \\[2pt]
\hline
\endhead
\hline
\endfoot
0 & Systemprozesse \\
1 & Interaktive Prozesse \\
2 & Allgemeine Prozesse \\
4 & Batch \\
\end{tabularx}
}


\begin{itemize}
  \item Pro Stufe eigene Warteschlange; Wiedereingliederung möglich
  \item Innerhalb Klasse: RR
  \item \textbf{Multi-Level-Feedback-Scheduling}: Prozess kann Klasse wechseln (häufigste Praxis-Lösung)
\end{itemize}


\section{Interrupts}

\textbf{Interrupt}: Unterbrechung der sequentiellen Programmausführung; Kontrolle an \textbf{ISR}. Synchron/asynchron.


\textbf{Befehlszyklus}: Fetch → Execute → Prüfung IRQ? → ggf. PC auf ISR-Adresse.


\textbf{Auslöser}: Speicherschutzverletzung, Hardware-Taktgeber, E/A-Gerät, Systemcalls etc.



{\small
\begin{tabularx}{\linewidth}{XXX}
\hline
\rowcolor{tableheadbg}
{\color{white}\textbf{Verfahren}} & {\color{white}\textbf{Vorteil}} & {\color{white}\textbf{Nachteil}} \\[2pt]
\hline
\endhead
\hline
\endfoot
Polling & leicht & CPU-Verschwendung \\
Interrupt & effizient & komplexer \\
\end{tabularx}
}



\subsection{Interruptklassen}


{\small
\begin{tabularx}{\linewidth}{XXX}
\hline
\rowcolor{tableheadbg}
{\color{white}\textbf{Klasse}} & {\color{white}\textbf{Eigenschaft}} & {\color{white}\textbf{Beispiele}} \\[2pt]
\hline
\endhead
\hline
\endfoot
Asynchron & nicht vorhersehbar/reproduzierbar & Hardware-Interrupts (Plattenblock, Netzwerk) \\
Synchron (Exception) & vorhersehbar/reproduzierbar & Systemcalls, Traps, Faults \\
\end{tabularx}
}


Synchron-Unterklassen:

\begin{itemize}
  \item \textbf{Fault}: Unterbrechung \textbf{vor} Ausführung (z.B. Page Fault)
  \item \textbf{Trap}: Unterbrechung \textbf{nach} Ausführung (z.B. Division durch 0)
\end{itemize}


\subsection{Bearbeitung}

Komponenten: Interrupt-Controller (IC), ISR, IRQ-Bearbeitung.


\textbf{Interrupt Controller}: nimmt Signale entgegen, informiert CPU.


\textbf{ISR (Interrupt Service Routine)}: einem Interrupt zugeordnete Anweisungsfolge; pro Interrupt-Typ eine ISR; BS stellt ISR für alle bereit; bei externen Geräten im Treiber.


\textbf{Benötigt}:

\begin{enumerate}
  \item \textbf{Interrupt-Vektor-Tabelle (IVT)}: an vordefinierter Stelle im Kernelspeicher; Eintrag = ISR-Adresse; Index = Interrupt-Nummer; IC bildet IRQ → Interrupt-Nummer ab
  \item \textbf{Interrupt-Level (IL)}: Priorität pro Interrupt
\end{enumerate}

\textbf{Ablauf}:

\begin{enumerate}
  \item Interrupt unterbricht Programm (wenn IL erlaubt)
  \item Prozessorstatus sichern (Stack)
  \item ISR via IVT suchen + ausführen (Kernelmodus)
  \item Bestätigung an IC
  \item Status wiederherstellen, Programm fortsetzen
\end{enumerate}

\begin{quote}
\textbf{Merke}: Ergebnis darf nicht von Interrupt-Auftreten abhängen. Bei gleichzeitigen IRQs entscheidet Priorität; gleich-/niedriger-priorisierte werden während Bearbeitung nicht angenommen.
\end{quote}


\subsection{Maskierung}
\begin{itemize}
  \item \textbf{IMR (Interrupt Mask Register)}: Bit = 1 → Interrupt aus
  \item IC wiederholt Anfrage bis angenommen
\end{itemize}


{\small
\begin{tabularx}{\linewidth}{XX}
\hline
\rowcolor{tableheadbg}
{\color{white}\textbf{Typ}} & {\color{white}\textbf{Bedeutung}} \\[2pt]
\hline
\endhead
\hline
\endfoot
MI & Maskable Interrupts \\
NMI & Non-Maskable Interrupts \\
Präzise & Befehle vollständig oder gar nicht ausgeführt \\
Unpräzise & teilweise Ausführung möglich \\
\end{tabularx}
}



\section{Synchronisation}

\textbf{Nebenläufigkeit}: quasi-parallele Ausführung von Prozessen/Threads.

\begin{itemize}
  \item 1 CPU/Kern → quasi-parallel
  \item mehrere → echt-parallel
\end{itemize}

\textbf{Synchronisation}: Mechanismen gegen Probleme bei nebenläufigem Zugriff auf gemeinsame Betriebsmittel.


\textbf{Probleme}: Blockieren, Verhungern, Verklemmung.



\subsection{Race Condition}

\textbf{Race Condition}: Ergebnis hängt von zeitlicher Reihenfolge konkurrierender Zugriffe ab. Entsteht durch Nebenläufigkeit \textbf{+} gemeinsame Betriebsmittel.


\textbf{Lost Update}: A liest 0, B liest 0, A schreibt, B schreibt → A's Update verloren.



\subsection{Kritischer Abschnitt}

\textbf{Kritischer Abschnitt}: Programmteil, der nur von einem Prozess gleichzeitig durchlaufen werden darf.


\textbf{Atomare Aktion}: logisch nicht unterbrechbarer Codebereich.


\textbf{Bedingungen (Mutex)}:

\begin{enumerate}
  \item \textbf{Mutual Exclusion}: nie zwei gleichzeitig drin
  \item Keine Annahmen über Geschwindigkeit/Anzahl Prozessoren
  \item Außerhalb darf kein Prozess andere blockieren
  \item \textbf{Fairness}: jeder Wartende muss irgendwann eintreten
\end{enumerate}

\textbf{Implementierung}:

\begin{enumerate}
  \item Busy Waiting — ständige Sperrabfrage; CPU-Verschwendung
  \item Schlafen + Aufwecken
\end{enumerate}


\subsection{Synchronisationskonzepte}

\begin{enumerate}
  \item Software (selten in BS)
  \item Interrupts sperren — alle Interrupts maskieren; Multiprozessor unpraktisch
  \item Hardware-Befehle (TSL/TAS)
\end{enumerate}


\subsubsection{Sperren (Lock)}
\begin{itemize}
  \item Sperrvariable: 0 = frei, 1 = belegt
  \item Operationen: Testen + Setzen
  \item \textbf{Spinlock}: Warteschleife (kurze Wartezeiten)
\end{itemize}

\textbf{TSL/TAS (Test and Set Lock)}: Test+Set ununterbrechbar in einem Maschinenbefehl; ein Speicherzugriffszyklus; CPU blockiert Bus; Multiprocessing-tauglich.



\subsubsection{Semaphor}

\textbf{Semaphor}: Sperrmechanismus mit Zähler + Warteschlange.



{\small
\begin{tabularx}{\linewidth}{XX}
\hline
\rowcolor{tableheadbg}
{\color{white}\textbf{Op}} & {\color{white}\textbf{Wirkung}} \\[2pt]
\hline
\endhead
\hline
\endfoot
\textbf{P()} / Down (anfordern) & Zähler > 0: Zähler−−; sonst: in Warteschlange + suspendiert \\
\textbf{V()} / Up (freigeben) & Zähler++; ggf. Wartenden aufwecken \\
\end{tabularx}
}


\begin{itemize}
  \item \textbf{Mutex / binäres Semaphor}: Zähler nur 0/1
  \item \textbf{Zählsemaphor}: beliebiger Initialwert → max. n gleichzeitig im KA
\end{itemize}

Nutzung: \texttt{P(); KA; V();}

\begin{itemize}
  \item \texttt{V();...P();} → alle drin
  \item \texttt{P();...P();} → niemand mehr drin
\end{itemize}

\begin{quote}
Semaphore sind fehleranfällig.
\end{quote}


\subsubsection{Monitor}

\textbf{Monitor}: Objekt mit gemeinsamen Daten, Zugriffsmethoden + „Türsteher" — nur 1 Prozess gleichzeitig aktiv. Compiler erzeugt Synchronisation.


\begin{itemize}
  \item \textbf{Bedingungsvariablen} (Conditions)
  \item Operationen: \textbf{Wait} (in Wartezustand), \textbf{Signal} (Wartenden aufwecken)
\end{itemize}

Ablauf:

\begin{enumerate}
  \item Bedingung nicht erfüllt → \texttt{wait(cond)} → Monitor temporär verlassen → anderer kann eintreten
  \item Anderer Prozess ändert Bedingung → \texttt{signal(cond)} → Wartender setzt fort
\end{enumerate}

Bsp. (Puffer): voll → \texttt{wait(full)}; entnommen → \texttt{signal(full)}.



{\small
\begin{tabularx}{\linewidth}{XXX}
\hline
\rowcolor{tableheadbg}
{\color{white}\textbf{Konzept}} & {\color{white}\textbf{Typ}} & {\color{white}\textbf{Verhalten}} \\[2pt]
\hline
\endhead
\hline
\endfoot
Signal-and-Continue & Mesa & Signalisierender arbeitet weiter \\
Signal-and-Wait & Hoare & Signalisierender verlässt Monitor sofort \\
\end{tabularx}
}



\section{Deadlock}

\textbf{Deadlock (Verklemmung)}: Mehrere Prozesse halten Betriebsmittel und warten auf weitere, die andere halten → Zyklus.



{\small
\begin{tabularx}{\linewidth}{XX}
\hline
\rowcolor{tableheadbg}
{\color{white}\textbf{Begriff}} & {\color{white}\textbf{Definition}} \\[2pt]
\hline
\endhead
\hline
\endfoot
Blockieren & P2 wartet auf von P1 gehaltene Ressource \\
Verhungern (Starvation) & Prozess erhält trotz Bereitschaft keine CPU-Zeit \\
\end{tabularx}
}



\subsection{Vier Bedingungen (alle gleichzeitig)}
\begin{enumerate}
  \item \textbf{Mutual Exclusion}: Ressource exklusiv
  \item \textbf{Hold and Wait}: hält und fordert weitere an
  \item \textbf{No Preemption}: kein Entzug
  \item \textbf{Circular Wait}: zyklische Warteabhängigkeit
\end{enumerate}


\subsection{Strategien}
\begin{enumerate}
  \item \textbf{Erkennen}: Betriebsmittelgraph (Knoten P/R, Kanten = Belegung/Anforderung), Zyklussuche
  \item \textbf{Ignorieren} (Vogel-Strauß): wartet unendlich; extern abbrechen
  \item \textbf{Vermeiden}: BS so bauen, dass Deadlocks unmöglich
  \item \textbf{Verhindern}: Anforderung nur, wenn kein zukünftiger Deadlock; nach und nach anfordern
\end{enumerate}

\begin{quote}
Erkennung/Vermeidung/Verhinderung sehr aufwendig — vollständige Abdeckung kaum möglich.
\end{quote}


\subsection{Klassische Probleme}
\begin{itemize}
  \item \textbf{Producer-Consumer (Bounded Buffer)}: synchronisierter Zugriff auf begrenzten Puffer
  \item \textbf{Reader-Writer}: Schreiber + Leser auf Datenstruktur
  \item \textbf{Dining Philosophers}: 5 Philosophen, 5 Gabeln; je 2 Gabeln nötig; max 2 essen gleichzeitig; niemand verhungert
\end{itemize}


\section{Speicherverwaltung — Grundlagen}

\textbf{Realer Adressraum}: physisch verbauter Speicher.


\textbf{Physikalische Speicheradresse}: Adresse im physischen Speicher.


\textbf{Adressraum}: Menge möglicher Adressen (z. B. 4 GB → \{0…2³²−1\}).


\textbf{Speicherinhalte}:

\begin{itemize}
  \item Statisch: Code, Konstanten, (un)initialisierte Daten — direkt aus exe geladen
  \item Dynamisch: Heap, Stack — zur Laufzeit
\end{itemize}

\textbf{Layout}: Code/Konstanten unten, Heap wächst nach oben, Stack wächst von oben nach unten.


\textbf{Techniken}:

\begin{enumerate}
  \item Monoprogramming
  \item Multiprogramming mit festen Partitionen
  \item Multiprogramming mit variablen Partitionen + Swapping
  \item Multiprogramming + virtueller Speicher + Paging
\end{enumerate}


\section{Virtueller Speicher}

\textbf{Virtueller Adressraum}: Speicher, den ein Prozess vom BS bekommt — nur virtuell.

\begin{itemize}
  \item Jeder Prozess sieht eigenen Adressraum
  \item Aktuell benötigte Teile im RAM, Rest auf Sekundärspeicher
  \item Speicherbedarf darf > RAM sein
\end{itemize}

\textbf{Vorteile}:

\begin{enumerate}
  \item Programme > RAM möglich
  \item Mehr Programme im RAM → CPU-Auslastung ↑
  \item Ansprechbarer Speicher > physisch
  \item Externe Fragmentierung untergeordnet
\end{enumerate}

\textbf{Seite (Page)}: Block des virtuellen Speichers.


\textbf{Seitenrahmen (Page Frame)}: Block des RAM. Alle gleich groß (typisch 1, 4, 8, 16, 64 KiB).


\textbf{Paging Area (Schattenspeicher)}: Festplattenbereich für ausgelagerte Seiten.



{\small
\begin{tabularx}{\linewidth}{XXX}
\hline
\rowcolor{tableheadbg}
{\color{white}\textbf{}} & {\color{white}\textbf{Swapping}} & {\color{white}\textbf{Paging}} \\[2pt]
\hline
\endhead
\hline
\endfoot
Granularität & ganzer Prozess & Teile \\
Einsatz & ohne MMU & mit MMU \\
\end{tabularx}
}



\section{MMU (Memory Management Unit)}

\textbf{MMU}: Funktionseinheit der CPU zur Adressübersetzung; ersetzt Basis-/Limitregister.

\begin{itemize}
  \item Bildet virtuelle → reale Adressen ab
  \item Pfad: \texttt{CPU → (virt.) → MMU → (real) → Hauptspeicher}
  \item Greift auf Seitentabellen zu
\end{itemize}


\section{Seitentabellen}

\textbf{Zweck}: Umrechnung virtuell → physisch. Pro virtuellem Adressraum eine Tabelle. Enthält Verweise auf Frames, \textbf{nicht} Daten.


\textbf{Adressaufbau}: virt. Adresse = Index \texttt{i} | Offset \texttt{d} → real. Adresse = Frame \texttt{f} | \texttt{d}.


\textbf{Seitentabelleneintrag}:


{\small
\begin{tabularx}{\linewidth}{XX}
\hline
\rowcolor{tableheadbg}
{\color{white}\textbf{Feld}} & {\color{white}\textbf{Bedeutung}} \\[2pt]
\hline
\endhead
\hline
\endfoot
P/A & Present/Absent (im Frame?) \\
Protection (C/R/M) & r/w/x \\
Frame-Number f & Verweis \\
M (Modified/Dirty) & verändert? \\
R (Reference) & Zugriff erfolgt? \\
C & Caching-Bit \\
\end{tabularx}
}



\subsection{Arten}


\subsubsection{Einstufig}
\begin{itemize}
  \item Direkter Index → Frame
  \item Bsp. 16-Bit virt. → 15-Bit real (16 Einträge bei 4 KiB Pages)
  \item 32-Bit / 4 KiB → >1.000.000 Einträge
  \item \textbf{Nachteil}: enormer Speicherbedarf
\end{itemize}


\subsubsection{Mehrstufig}
\begin{itemize}
  \item Adresse: \texttt{i1 | i2 | d}
  \item 32-Bit: 10+10+12 → 1024 × 1024 Pages à 4 KiB
  \item Top-Level + Second-Level
  \item \textbf{Vorteil}: Subtabellen müssen nicht alle im Speicher
  \item \textbf{Nachteil}: häufige, aufwändige Berechnungen
\end{itemize}


\subsubsection{Invertiert}
\begin{itemize}
  \item Eintrag pro physischem Frame (nicht pro virtueller Seite)
  \item Eintrag: \texttt{pid | i | f} + Kollisionsliste
  \item Zugriff via 12-Bit Hash auf \texttt{(pid, i)}
  \item \textbf{Vorteil}: viel weniger Einträge; komplett im RAM
  \item \textbf{Nachteil}: keine Ordnung nach virtuellen Adressen, Suche aufwändig
\end{itemize}


\subsubsection{TLB (Translation Lookaside Buffer)}
\begin{itemize}
  \item Schneller Cache in CPU für aktive (virt → frame) Mappings
  \item \textbf{TLB hit}: direkt
  \item \textbf{TLB miss}: Tabellenlookup
\end{itemize}


\section{Page Fault \& Seitenersetzung}

\textbf{Page Fault}: Trap, wenn benötigte virtuelle Seite nicht im Frame eingelagert ist. Bearbeitung im Kernelmodus.



{\small
\begin{tabularx}{\linewidth}{XX}
\hline
\rowcolor{tableheadbg}
{\color{white}\textbf{Begriff}} & {\color{white}\textbf{Definition}} \\[2pt]
\hline
\endhead
\hline
\endfoot
Eingelagerte Seite & Inhalt liegt in Frame \\
Ausgelagerte Seite & nicht in Frame \\
Modifizierte Seite & seit Einlagerung verändert (M-Bit=1) \\
\end{tabularx}
}


\begin{quote}
Modifizierte Seiten dürfen nicht ohne Sicherung ersetzt werden.
\end{quote}


{\small
\begin{tabularx}{\linewidth}{XXX}
\hline
\rowcolor{tableheadbg}
{\color{white}\textbf{}} & {\color{white}\textbf{Demand Paging}} & {\color{white}\textbf{Prepaging}} \\[2pt]
\hline
\endhead
\hline
\endfoot
Politik & nur auf Anforderung & spekulativ vor Zugriff \\
Vorteil & nur Benötigtes geladen & nutzt Lokalität \\
\end{tabularx}
}



\subsection{Ablauf Seitenersetzung}
\begin{enumerate}
  \item MMU: Seite B nicht eingelagert → Page Fault
  \item Kein freier Frame → SEA (Seitenersetzungsalgorithmus) startet
  \item Opfer-Seite A in Frame X bestimmen
  \item Wenn A modifiziert: nach Paging Area schreiben
  \item Seite B in Frame X einlagern
\end{enumerate}


\section{Seitenersetzungsalgorithmen (SEA)}


\subsection{Optimal (Belady)}
\textbf{Optimal}: Verdrängt Seite, die am längsten in Zukunft nicht gebraucht wird.

\begin{itemize}
  \item Praktisch unmöglich (Zukunft unbekannt) → nur Referenz
\end{itemize}


\subsection{FIFO}
\textbf{FIFO}: Verdrängt älteste eingelagerte Seite.

\begin{enumerate}
  \item Liste der Frames absteigend nach Einlagerungszeit
  \item Page Fault: ältester (vorne) verdrängt; neue Seite ans Ende
  \item R-Bit nicht benötigt
  \item Nimmt keine Rücksicht auf Nutzung
\end{enumerate}


\subsection{NRU (Not Recently Used)}
\textbf{NRU}: 4 Klassen aus (R, M); R-Bits periodisch zurücksetzen (Tanenbaum: \textasciitilde{}20 ms). Klasse 1 zuerst auslagern.



{\small
\begin{tabularx}{\linewidth}{XXXX}
\hline
\rowcolor{tableheadbg}
{\color{white}\textbf{Klasse}} & {\color{white}\textbf{R}} & {\color{white}\textbf{M}} & {\color{white}\textbf{Verhalten}} \\[2pt]
\hline
\endhead
\hline
\endfoot
1 & 0 & 0 & nicht ref., nicht mod. — keine Sicherung nötig \\
2 & 0 & 1 & nicht ref., mod. — schreiben \\
3 & 1 & 0 & ref., nicht mod. — wahrscheinl. wieder gebraucht \\
4 & 1 & 1 & ref., mod. — schreiben + wieder nötig \\
\end{tabularx}
}



\subsection{Second Chance}
\textbf{Second Chance}: Verbessert FIFO durch R-Bit-Prüfung der ältesten Seite.

\begin{itemize}
  \item R=1 → R:=0, ans Ende; R=0 → auslagern
  \item Degeneriert zu FIFO, wenn alle R=1
\end{itemize}


\subsection{Clock Page}
\textbf{Clock Page}: Second Chance als zirkuläre Liste mit Zeiger.

\begin{itemize}
  \item R=0 → auslagern; R=1 → R:=0, Zeiger++
  \item Max. 1 Runde
\end{itemize}


\subsection{LRU (Least Recently Used)}
\textbf{LRU}: Verdrängt am längsten ungenutzte Seite (Annahme: kürzlich genutzt → bald wieder).

\begin{itemize}
  \item \textbf{Matrix-Verfahren} (n×n): Zugriff auf k → Zeile k=1, Spalte k=0; Max=zuletzt; Min=Opfer
  \item Alternative: Zeitstempel pro Eintrag
  \item Nachteil: hoher Aufwand pro Zugriff
  \item Praxis: Pseudo-LRU (Clock, Second Chance) mit R/M
\end{itemize}


\subsection{NFU (Not Frequently Used)}
\textbf{NFU}: Zähler pro Seite, Auslagerung des kleinsten Zählers.

\begin{itemize}
  \item Bei Zugriff: R-Bit setzen
  \item Timer: Zähler += R; R := 0
  \item Problem: alte hohe Zähler werden nie verdrängt
\end{itemize}


\subsection{Aging}
\textbf{Aging}: Verbessert NFU gegen Aufholproblem durch Bit-Shift statt Addition.

\begin{itemize}
  \item Vor Addition: Zähler 1 Bit nach rechts schieben
  \item R-Bit wird höchstwertiges Bit
  \item ≈ klassisches LRU
\end{itemize}


\subsection{Working Set}
\textbf{Working Set}: Aktuell benötigte Seitenmenge; ständig im RAM → wenige Page Faults.

\begin{itemize}
  \item Verdrängt: Seiten außerhalb WS
  \item \textbf{Lokalitätseffekt} (80/20): nur kleiner Code-Teil aktiv
  \item Distanz Δ = wichtiger Parameter
\end{itemize}


\section{Lokalität}

\textbf{Lokalität}: Prinzip, dass Speicherzugriffe nicht zufällig verteilt sind, sondern Muster aufweisen — genutzt zur Cache-Optimierung.



{\small
\begin{tabularx}{\linewidth}{XXX}
\hline
\rowcolor{tableheadbg}
{\color{white}\textbf{}} & {\color{white}\textbf{Zeitlich}} & {\color{white}\textbf{Räumlich}} \\[2pt]
\hline
\endhead
\hline
\endfoot
Bezug & zeitlicher Ablauf & räumlicher Abstand \\
Aussage & Adresse bald wieder zugegriffen & Nachbaradresse bald genutzt \\
Optimierung & im Cache halten & Nachbarn mitladen \\
\end{tabularx}
}



\section{Fragmentierung}


{\small
\begin{tabularx}{\linewidth}{XXX}
\hline
\rowcolor{tableheadbg}
{\color{white}\textbf{Art}} & {\color{white}\textbf{Wo}} & {\color{white}\textbf{Bedeutung}} \\[2pt]
\hline
\endhead
\hline
\endfoot
Intern & innerhalb belegter Blöcke & Block > Bedarf, Rest verschwendet (Ø 50\% letzte Seite) \\
Extern & zwischen belegten Blöcken & Lücken einzeln zu klein, summiert ausreichend \\
\end{tabularx}
}


\textbf{Speicherarten}:

\begin{itemize}
  \item RAM: beide Arten
  \item HDDs: beide Arten
  \item SSDs: vor allem intern; extern weniger problematisch
\end{itemize}

\textbf{Defragmentierung}: Reorganisation belegter Blöcke zu zusammenhängenden Bereichen → bessere Performance (v.a. HDDs).



\section{Placement (Speicherbelegung)}

\textbf{Ziele}: Fragmentierung minimieren, wenig Speicher, schnell Frame finden.


\textbf{Verwaltung freier Bereiche}:

\begin{enumerate}
  \item \textbf{Bitmap}: 1 Bit/Frame (1=belegt, 0=frei); freie Bereiche = Folge von 0; weniger Speicher als Liste
  \item \textbf{Verkettete Liste}: freie Bereiche
\end{enumerate}


{\small
\begin{tabularx}{\linewidth}{XX}
\hline
\rowcolor{tableheadbg}
{\color{white}\textbf{Strategie}} & {\color{white}\textbf{Vorgehen}} \\[2pt]
\hline
\endhead
\hline
\endfoot
First-Fit & erster ausreichend großer Bereich \\
Next-Fit & wie First, ab letzter Position \\
Best-Fit & kleinste passende Lücke \\
Worst-Fit & größte Lücke \\
Buddy-Technik & 2ᵏ-Halbierung \\
Quick-Fit & wie Buddy, aber nicht doppelte Größe \\
\end{tabularx}
}



\subsection{Buddy-Technik}
\begin{itemize}
  \item Blockgrößen = 2ᵏ; Anforderung auf nächste 2ᵏ aufgerundet
  \item Halbierung bei Anforderung; benachbarte gleichgroße freie = Buddies
  \item Bei Freigabe: Buddies verschmolzen
  \item \textbf{Vorteil}: passende Größe oft verfügbar; gegen externe Fragmentierung
  \item \textbf{Nachteil}: zunehmende interne Fragmentierung durch Aufrundung
\end{itemize}


\subsection{Quick-Fit}
\begin{itemize}
  \item Ähnlich Buddy, aber nächstes Level ≠ doppelte Größe
  \item Verschmelzen aufwändiger
\end{itemize}


\section{Festplatten und Cluster}

\textbf{Cluster}: Zusammenschluss benachbarter Sektoren, kleinste Adressierungseinheit; jede Datei belegt mind. 1 Cluster.

\begin{itemize}
  \item Reduziert Adressierungslast
  \item \textbf{Nachteil}: Fragmentierung möglich
\end{itemize}

\textbf{Slack-Arten}:


{\small
\begin{tabularx}{\linewidth}{XX}
\hline
\rowcolor{tableheadbg}
{\color{white}\textbf{Art}} & {\color{white}\textbf{Bedeutung}} \\[2pt]
\hline
\endhead
\hline
\endfoot
File Slack & ungenutzter Platz im letzten Cluster \\
RAM Slack & Bereich von EOF bis Sektorende, mit RAM-Inhalten gefüllt \\
Drive Slack & restliche Sektoren bis Clustergrenze, mit HDD-Daten gefüllt \\
\end{tabularx}
}


\begin{itemize}
  \item Linux: Zufallszahlen statt RAM/Drive-Inhalte
  \item Forensik: Slack-Inhalte wertvoll
\end{itemize}

\textbf{Cluster-Größe}:

\begin{itemize}
  \item Große Partition → große Cluster → Slack ↑ → ineffizient
  \item Kleine Partition → kleine Cluster → bessere Anpassung
\end{itemize}


\section{Geräteverwaltung}

\textbf{Geräteverwaltung}: Software für optimierte Zusammenarbeit BS↔Peripherie; Schnittstelle physische Geräte ↔ BS.

\begin{itemize}
  \item Prozesse kommunizieren \textbf{nicht direkt} mit Geräten → nur via Syscalls
  \item \textbf{Aufgaben}: Kommandos an Geräte, Interrupt-Bearbeitung, Fehlerbehandlung
\end{itemize}


\subsection{Schichtenarchitektur}
\begin{itemize}
  \item E/A-Geräte über Bussysteme an CPU
  \item \textbf{Northbridge}: interne Geräte (RAM, AGP)
  \item \textbf{Southbridge}: externe (PCI, S-ATA, SCSI, USB)
\end{itemize}

\textbf{Transferraten} (Auswahl):


{\small
\begin{tabularx}{\linewidth}{XX}
\hline
\rowcolor{tableheadbg}
{\color{white}\textbf{Gerät}} & {\color{white}\textbf{Rate}} \\[2pt]
\hline
\endhead
\hline
\endfoot
Keyboard & \textasciitilde{}10 B/s \\
Maus & \textasciitilde{}100 B/s \\
Laserdrucker & \textasciitilde{}100 KB/s \\
ISA-Bus & 16 MB/s \\
USB1/2/3 & 12 Mbit/s / 480 Mbit/s / 5 Gbit/s \\
Fast/Gigabit Ethernet & 100 Mbit/s / bis 10 Gbit/s \\
PCIe & 32 Gbit/s \\
ATA / SATA & 130 MB/s / \textasciitilde{}2 GB/s \\
\end{tabularx}
}



\subsection{Gerätetreiber}

\textbf{Gerätetreiber}: Softwaremodul, das Interaktionen zwischen BS und Controller eines Geräts steuert.


\textbf{Aufgaben}:

\begin{itemize}
  \item Gerät dem BS bekannt machen
  \item Controller initialisieren
  \item Befehle gerätespezifisch übersetzen
  \item Datenpufferung Gerät ↔ Hauptspeicher
  \item ISR (Unterbrechungsbearbeitung) enthalten
  \item Koordination nebenläufiger Zugriffe
\end{itemize}


\subsection{Geräteklassen}


\subsubsection{Blockorientiert}
\begin{itemize}
  \item Übertragung in \textbf{kompletten Blöcken}, direkt adressierbar
  \item Blockgrößen 512–32.768 Byte
  \item Bsp.: Festplatte, Bandlaufwerk, CD/DVD
  \item Funktionen: \texttt{initDevice}, \texttt{readBlock}, \texttt{writeBlock}, \texttt{handleInterrupt}
\end{itemize}


\subsubsection{Zeichenorientiert}
\begin{itemize}
  \item \textbf{Datenstrom (stream)} aus Zeichen
  \item Bsp.: Maus, Tastatur, Netzwerkkarte, serielle/parallele/USB
  \item Funktionen: \texttt{initDevice}, \texttt{readChar}, \texttt{writeChar}, \texttt{handleInterrupt}
\end{itemize}


\subsection{Controller (Gerätesteuereinheit)}
\begin{itemize}
  \item Kommunikation Hardware ↔ BS
  \item \textbf{Register} zur CPU-Kommunikation: Schreiben → Befehle ans Gerät; Lesen → Status
  \item Anwendungen via Syscalls: \texttt{open}, \texttt{close}, \texttt{read}, \texttt{write}
\end{itemize}


\subsection{Adressierung Controller-Register}


\subsubsection{I/O-Port}
\begin{itemize}
  \item \textbf{Zwei Adressräume}: Speicher und I/O getrennt
  \item Portnummer (8/16-Bit) pro Register, geschützter Bereich
\end{itemize}


\subsubsection{Memory-Mapped I/O}
\begin{itemize}
  \item \textbf{Ein Adressraum}: I/O + Speicher gemeinsam
  \item E/A-Adresse pro Gerät, meist im oberen Adressraum
  \item Adressen dürfen nicht kollidieren → automatische Vergabe
\end{itemize}


{\small
\begin{tabularx}{\linewidth}{XXX}
\hline
\rowcolor{tableheadbg}
{\color{white}\textbf{Methode}} & {\color{white}\textbf{Vorteil}} & {\color{white}\textbf{Nachteil}} \\[2pt]
\hline
\endhead
\hline
\endfoot
I/O-Port & einfach & spezielle Port-Befehle nötig \\
MMIO & Zugriff wie RAM, leichter Schutz & Unterscheidung echt/gemappt aufwändig \\
\end{tabularx}
}


\begin{quote}
CPU-Adressierung der Register kostet bei beiden Methoden sehr viel Rechenzeit.
\end{quote}


\subsection{DMA (Direct Memory Access)}

\textbf{DMA}: Transport ganzer Datenblöcke direkt Geräte-Controller ↔ Hauptspeicher via DMA-Controller; CPU nur Anfang/Ende beteiligt.


\textbf{Vorteile}:

\begin{itemize}
  \item CPU-Entlastung
  \item Höhere Gesamtgeschwindigkeit
  \item Weniger Interrupts (1 statt 1 pro Wort)
  \item DMA-Controller spezialisiert
\end{itemize}

\textbf{Ablauf (Lesen von Platte)}:

\begin{enumerate}
  \item CPU programmiert DMA-Controller (Quelle, Ziel, Anzahl Wörter)
  \item CPU sendet Lese-Kommando an Platten-Controller
  \item DMA fordert Platten-Controller zur Übertragung an
  \item Platte → DMA → RAM (über Bus)
  \item Platten-Controller meldet DMA Abschluss
  \item DMA erzeugt \textbf{einen} Interrupt → CPU
  \item ISR ausgeführt
\end{enumerate}

\begin{itemize}
  \item DMA-Zugriffe laufen parallel zur CPU-Bearbeitung
\end{itemize}


\section{Dateiverwaltung}

\textbf{Dateiverwaltung}: organisiert dauerhafte Daten (HDD/SSD); nutzt Geräteverwaltung; stellt Dateidienste bereit.


\textbf{Operationen (CRUD)}: Create, Read, Update, Delete.


\textbf{Systemdienste}: Erzeugen / Löschen / Kopieren von Dateien und Katalogeinträgen, Verknüpfungen, Umbenennen.



\section{Dateisystem}

\textbf{Dateisystem}: ermöglicht geordnete Ablage + Wiederfinden auf Speichermedium; stellt Zugriff bereit.


\textbf{Aufgaben}: CRUD, Verwaltung freier/belegter Speicher, Strukturierung, Verzeichnisse, Attribute, Zugriffsrechte.


\textbf{Speicherung}:

\begin{itemize}
  \item Variante 1: n aufeinanderfolgende Bytes (Größenänderung → Verschiebung)
  \item Variante 2: Aufteilung in Blöcke (Standard, flexibel)
\end{itemize}

\textbf{Beispiele}:


{\small
\begin{tabularx}{\linewidth}{XX}
\hline
\rowcolor{tableheadbg}
{\color{white}\textbf{BS}} & {\color{white}\textbf{Dateisysteme}} \\[2pt]
\hline
\endhead
\hline
\endfoot
Apple & Apple DOS, ProDOS, MFS, HFS/HFSX/HFS+, APFS \\
Linux & ext, ext2, ext3, ext4, ReiserFS \\
Windows & FAT12/16/32, exFAT, VFAT, NTFS, ReFS \\
\end{tabularx}
}



\section{Dateisystemcheck}

\textbf{Grundproblem}: nach Absturz Konsistenz herstellen. Optionen:

\begin{enumerate}
  \item Strukturen immer konsistent halten
  \item Veränderungen sequenziell ins Journal/Log
\end{enumerate}


\subsection{Journaling-Dateisystem}

\textbf{Prinzip}: Jeder Vorgang einer Transaktion: erst Log-Eintrag, \textbf{dann} Änderung. Log temporär, Daten später ans Ziel zurückgeschrieben.


\begin{quote}
Log-Eintrag steht \textbf{immer vor} der Änderung auf Platte.
\end{quote}

\textbf{Transaktion}: atomar; nur 2 Zustände — komplett oder nicht.

\begin{itemize}
  \item Nicht beendete: rückgängig
  \item Beispiele: Erzeugen/Löschen/Erweitern/Verkürzen/Umbenennen, Attribute ändern
\end{itemize}

\textbf{Inkonsistenzen} entstehen bei Stromausfall, HW-Ausfall, Absturz.

\begin{itemize}
  \item Bsp.: Verzeichniseintrag fehlt, Block belegt aber nicht markiert
\end{itemize}

\textbf{Wiederherstellung}: beim Boot Protokoll abgleichen.



{\small
\begin{tabularx}{\linewidth}{XXX}
\hline
\rowcolor{tableheadbg}
{\color{white}\textbf{Art}} & {\color{white}\textbf{Inhalt}} & {\color{white}\textbf{Eigenschaft}} \\[2pt]
\hline
\endhead
\hline
\endfoot
Meta-Daten-Journaling & nur Operationen, kein Inhalt & Daten beim Absturz können verloren gehen \\
Vollständiges/Data-Journaling & gesamter Vorgang inkl. Daten & Original beim Überschreiben gesichert; doppeltes Schreiben \\
\end{tabularx}
}


Meta-Daten-Journaling protokolliert: Verzeichnisreferenz, Rechte, Eigentümer, Blöcke.



\subsection{Log-structured-Dateisystem}

\textbf{Prinzip}: Ganze Platte als Log; kein Zurückschreiben — Log \textbf{ist} die Organisation.

\begin{itemize}
  \item Alte Daten nicht überschreiben, sondern anhängen (\textbf{Copy-on-Write})
  \item Verweise auf modifizierte Kopien umgebogen
\end{itemize}

\textbf{Aufbau}: Log Head → Daten → Empty Space → Daten → Log Tail.


\textbf{Schreiben}:

\begin{enumerate}
  \item Nicht bei jeder Änderung
  \item Schreibaufträge im RAM gepuffert
  \item Periodisch als Segment ans Ende des Logs
\end{enumerate}

\textbf{Cleaner}:

\begin{enumerate}
  \item Lesen Zusammenfassung erstes Segment (welche I-Nodes/Dateien)
  \item Prüfung I-Node-Map: noch aktuell?
\begin{itemize}
  \item JA: in Speicher laden, mit nächstem Segment neu schreiben
  \item NEIN: Info löschen, Segment freigeben
\end{itemize}
\end{enumerate}

\textbf{Nachteile}:

\begin{itemize}
  \item Protokollierungsaufwand (Platte springt zw. Log und Daten)
  \item Doppeltes Schreiben (Data-Journaling)
  \item Starke Fragmentierung (Log-structured)
\end{itemize}


\section{Datensicherung / Backup}

\textbf{Adressierte Probleme}: Katastrophe, Unachtsamkeit.



{\small
\begin{tabularx}{\linewidth}{XXXX}
\hline
\rowcolor{tableheadbg}
{\color{white}\textbf{Art}} & {\color{white}\textbf{Basis}} & {\color{white}\textbf{Speicher}} & {\color{white}\textbf{Wiederherstellung}} \\[2pt]
\hline
\endhead
\hline
\endfoot
\textbf{Vollbackup} & — & sehr hoch & nur Vollbackup \\
\textbf{Inkrementell} & letztes Backup & niedrig & Voll + alle Inkremente \\
\textbf{Differentiell} & letztes Vollbackup & mittel & Voll + gewünschtes Diff \\
\end{tabularx}
}


\begin{itemize}
  \item \textbf{Voll}: 1 Band, Wiederherstellung zu jedem Zeitpunkt
  \item \textbf{Inkrementell}: kleine Inkremente; fehlt eines, fehlen Dateien
  \item \textbf{Differentiell}: einmal veränderte Dateien werden bei jedem Diff erneut gesichert
\end{itemize}


\section{Meltdown \& Spectre}

\textbf{Meltdown}: Sicherheitslücke (3.1.2018), die geschützten RAM (Passwörter etc.) via CPU-Designfehler ausliest.

\begin{itemize}
  \item Entdeckt: Google Project Zero (28.7.2017)
  \item Betroffen: alle Intel-Chips ab 1995
  \item Hebelt RAM-Trennung aus via \textbf{Seitenkanalangriff} (Cache)
\end{itemize}


\subsection{Grundlagen}

\textbf{Kernel}: Herzstück des BS mit voller Kontrolle, vom Boot bis Speicherübergabe.



{\small
\begin{tabularx}{\linewidth}{XXX}
\hline
\rowcolor{tableheadbg}
{\color{white}\textbf{Modus}} & {\color{white}\textbf{Rechte}} & {\color{white}\textbf{Zugriff}} \\[2pt]
\hline
\endhead
\hline
\endfoot
Kernel & uneingeschränkt & gesamte HW + Speicher \\
User & eingeschränkt & nur via System-API \\
\end{tabularx}
}


\textbf{Speicherisolation}: pro Prozess eigener virtueller Adressraum + Page-Table; Kontextwechsel = Flush + Laden (teuer).


\textbf{Kernel-Einbettung}: Aus Performancegründen ist Kernel-Adressraum + gesamter physischer Speicher in \textbf{jeder} User-Page-Table eingebettet, geschützt durch \textbf{Supervisor-Bit}.

\begin{itemize}
  \item User-Zugriff auf Kernel-Adresse → \textbf{Seg-Fault}
\end{itemize}


\subsection{CPU-Optimierungen}


\subsubsection{Pipelining (5 Stages)}

{\small
\begin{tabularx}{\linewidth}{XX}
\hline
\rowcolor{tableheadbg}
{\color{white}\textbf{Stage}} & {\color{white}\textbf{Aufgabe}} \\[2pt]
\hline
\endhead
\hline
\endfoot
FETCH & Instruktion holen \\
DECODE & dekodieren \\
LOAD & laden \\
EXECUTE & ausführen \\
WRITE & Ergebnis schreiben \\
\end{tabularx}
}


Mehrere Instruktionen überlappen.



\subsubsection{Out-of-Order Execution}
\begin{itemize}
  \item Befehle nicht in Programmreihenfolge, sondern optimiert ausgeführt
  \item Wartet auf Daten → späteren Befehl vorziehen
\end{itemize}


\subsubsection{Speculative Execution}
\begin{itemize}
  \item CPU rät, welcher Code als nächstes kommt; führt vorab aus
  \item Fehlrate → Ergebnis verworfen
  \item \textbf{Aber}: Cache und TLB-Updates \textbf{nicht zurückgesetzt}
  \item Bei Spekulation gilt Schutzregel \textbf{nicht} → Exploit-Tür
\end{itemize}


\subsection{Meltdown-Angriff}

\textbf{Prinzip}: Kombination aus Speculative Execution + nicht-zurückgesetztem Cache + Cache-Timing.


\textbf{Ablauf} (alles im User-Mode):

\begin{enumerate}
  \item Speculative Execution greift unauthorisiert auf Kernel-Adresse zu
  \item Daten landen im Cache
  \item Exception (Seg-Fault) — Cache bleibt befüllt
  \item Cache-Timing-Attack (Flush \& Reload) durch zweiten Thread liest Geheimnis aus Latenzen
\end{enumerate}

\textbf{Probe-Array-Trick}:

\begin{lstlisting}
ProbeArray = [ASCII];
secret = readCharacter[5000];   // z.B. 80
character = ProbeArray[secret]; // ProbeArray[80] in Cache
\end{lstlisting}

Process 2 misst Zeiten — schnellster Index = Geheimnis. Iteration → Passwort rekonstruieren.



\subsection{Reparatur}

\textbf{Hardware-Abschaltung}: Out-of-Order/Branch Prediction/Speculative Execution abschalten = drastischer Performance-Verlust.


\textbf{OS-Patch}: \textasciitilde{}5–30\% Performance-Einbuße.


\textbf{KASLR (Kernel Address Space Layout Randomization)}: zufällige Kernel-Adressen zur Laufzeit.

\begin{itemize}
  \item Schwäche: Spraying und Side-Channels umgehen ASLR
\end{itemize}

\textbf{KAISER / KPTI (Kernel Page Table Isolation)}: zwei separate Adressräume (User vs. Kernel); Kernel nicht mehr in User-Page-Table eingebettet.

\begin{itemize}
  \item Verlangsamt Syscalls + Interrupts
\end{itemize}

\textbf{Zukunft}:

\begin{itemize}
  \item \textbf{Hardsplit}: Hardware-Aufteilung in feste Hälften
  \item Spekulatives Laden einschränken
  \item AMD: keine Spekulation auf geschützten Bereich
\end{itemize}

\begin{quote}
Spectre ist deutlich schwerer zu patchen.
\end{quote}


\subsection{Recap}
\begin{itemize}
  \item Designfehler in CPUs, kein SW-Bug
  \item Löschen/verschlüsseln keine Daten
  \item Nicht remote nutzbar — Code muss lokal laufen
  \item Hebeln Isolation Prozess↔Prozess + Prozess↔Kernel auf
  \item Meltdown: leichter auszunutzen + leichter zu beheben als Spectre
\end{itemize}




% ──────────────────────────────────────────────────────────────

\end{document}
